\chapter{立体几何初步}
\label{chap:space-geometry}

\tikzset{
  sg edge/.style={line width=.75pt,draw=BookBlue},
  sg hidden/.style={line width=.65pt,draw=BookBlue!60,dashed},
  sg aux/.style={line width=.6pt,draw=BookGold,dashed},
  sg plane/.style={draw=BookBlue,fill=BookBlueLight,fill opacity=.8},
  sg mark/.style={circle,fill=BookGold,inner sep=1.25pt},
  sg label/.style={font=\small,inner sep=1pt}
}
\newcommand{\sgcaption}[1]{\par\small\color{BookBlue!85!black}#1}

\section{点、直线、平面}

平面是向四周无限延展的；图中常画成平行四边形，记作平面 $\alpha$，也可用其中不共线的三点记作平面 $ABC$。基本语言如下。

\begin{center}
  \begin{tabular}{llll}
    \toprule
    关系 & 记号 & 关系 & 记号 \\
    \midrule
    点 $A$ 在直线 $l$ 上 & $A\in l$ & 点 $A$ 不在 $l$ 上 & $A\notin l$ \\
    点 $A$ 在平面 $\alpha$ 内 & $A\in\alpha$ & 点 $A$ 不在 $\alpha$ 内 & $A\notin\alpha$ \\
    直线 $l$ 在平面 $\alpha$ 内 & $l\subset\alpha$ & 直线 $l$ 与 $\alpha$ 相交 & $l\cap\alpha=A$ \\
    两平面交于直线 $l$ & $\alpha\cap\beta=l$ & 两平面平行 & $\alpha\parallel\beta$ \\
    \bottomrule
  \end{tabular}
\end{center}

点与直线之间用 $\in$；直线与平面之间用 $\subset$。不要写成 $l\in\alpha$。

\begin{center}
\begin{tikzpicture}[scale=.82]
  \path[sg plane] (0,0) -- (4.2,0) -- (5.1,1.6) -- (.9,1.6) -- cycle;
  \draw[sg edge] (-.25,.55) -- (4.65,.55) node[right,sg label] {$l$};
  \node[sg mark,label={[sg label]above:$A$}] at (1.25,.55) {};
  \node[sg mark,label={[sg label]above:$B$}] at (3.15,.55) {};
  \node[sg label] at (4.2,1.28) {$\alpha$};
  \draw[sg edge] (6.1,0) -- (10.3,0) -- (11.2,1.6) -- (7,1.6) -- cycle;
  \draw[sg edge,fill=BookGold!8] (7.1,-.15) -- (10.35,.2) -- (10.15,2) -- (6.9,1.65) -- cycle;
  \draw[very thick,BookGold] (7.02,.75) -- (10.24,.9) node[right,sg label] {$m$};
  \node[sg label] at (10.35,1.35) {$\alpha$};
  \node[sg label] at (7.35,1.75) {$\beta$};
\end{tikzpicture}
\sgcaption{两点定线在面内；两平面的公共点全部落在交线 $m$ 上。}
\end{center}

\subsection{平面的三个基本性质}

\begin{lawbox}[点线面的基本性质]
  \begin{enumerate}
    \item 若直线 $l$ 上两点 $A,B$ 在平面 $\alpha$ 内，则 $l\subset\alpha$。
    \item 不在同一直线上的三点确定唯一平面。
    \item 若两个不重合平面有一个公共点 $P$，则它们相交于一条过 $P$ 的直线：$\alpha\cap\beta=l$，$P\in l$。
  \end{enumerate}
\end{lawbox}

第二条立即给出三种常用的“确定平面”方式：

\[
  \boxed{\text{一条直线和线外一点}}
  \qquad
  \boxed{\text{两条相交直线}}
  \qquad
  \boxed{\text{两条平行直线}}.
\]

证明一条直线在平面内，通常只证直线上两点在该平面内；证明若干点共面，通常先用上述三种方式确定一个平面，再逐点验证。证明空间点共线，常把这些点同时放入两个平面，利用“两平面的公共点都在交线”这一事实。

\begin{examplebox}[原稿例：交线法证明共线]
  在四面体 $ABCD$ 中，两条分别位于面 $ABD$、面 $BCD$ 内的直线交于点 $P$。若已知 $P$ 同时在这两个面内，证明 $P,B,D$ 共线。

  \textbf{最短解：}$P\in\text{平面 }ABD$ 且 $P\in\text{平面 }BCD$；而两平面的交线为 $BD$，故 $P\in BD$。因此 $P,B,D$ 共线。

  \textbf{速查：}共线题先找两个面；对目标点逐一建立“双重归属”。
\end{examplebox}

\section{空间两直线的位置关系}

两直线若在同一平面内，则相交或平行；若不同在任何一个平面内，则称为\textbf{异面直线}。因此空间两直线只有三种位置关系：
\[
  \text{相交}\quad/\quad\text{平行}\quad/\quad\text{异面}.
\]

\begin{lawbox}[异面直线的快速判定]
  若直线 $a$ 与平面 $\alpha$ 相交于 $A$，直线 $b\subset\alpha$ 且 $A\notin b$，则 $a,b$ 异面。
\end{lawbox}

证明只需反证：若 $a,b$ 共面于 $\beta$，则 $A\in a\subset\beta$ 且 $b\subset\alpha\cap\beta$；由于 $A\notin b$，两平面的交线不能同时容纳 $A$ 与 $b$，矛盾。

\subsection{异面直线所成的角与公垂线}

过空间任一点 $O$ 分别作 $a'\parallel a$、$b'\parallel b$，相交直线 $a',b'$ 所成的锐角或直角称为异面直线 $a,b$ 所成的角。平移所选点不改变该角，这是等角定理的直接应用：若两个角的两边分别平行且方向相同，则两角相等。

若所成角为 $90^\circ$，记作 $a\perp b$；这一定义同时适用于相交直线和异面直线。两条异面直线有且只有一条公垂线段，其长度称为两异面直线间的距离。

\begin{center}
\begin{tikzpicture}[scale=.9]
  \coordinate (A) at (0,0); \coordinate (B) at (3.2,0);
  \coordinate (D) at (1.05,1.05); \coordinate (C) at (4.25,1.05);
  \coordinate (A1) at (0,2.6); \coordinate (B1) at (3.2,2.6);
  \coordinate (D1) at (1.05,3.65); \coordinate (C1) at (4.25,3.65);
  \draw[sg edge] (A)--(B)--(C)--(C1)--(B1)--(A1)--cycle (A1)--(D1)--(C1) (B1)--(C1);
  \draw[sg hidden] (A)--(D)--(C) (D)--(D1) (D1)--(C1);
  \draw[very thick,BookGold] (A1)--(B);
  \draw[very thick,BookBlue] (B1)--(C);
  \foreach \p/\pos in {A/below left,B/below,D/left,C/right,A1/left,B1/right,D1/above,C1/above}
    \node[sg label,\pos] at (\p) {$\p$};
  \node[sg label,align=left] at (7.15,2.5) {把 $B_1C$ 平移为 $A_1D$：\\[2pt]
    $\angle(A_1B,B_1C)=\angle BA_1D$};
  \coordinate (X) at (6.25,1.35); \coordinate (Y) at (8.4,1.35); \coordinate (Z) at (7.65,3.05);
  \draw[sg edge,-{Latex[length=2mm]}] (5.1,1.35)--(X);
  \draw[sg edge] (X)--(Y);
  \draw[very thick,BookGold] (X)--(Z);
  \pic[draw=BookGold,angle radius=6mm,"$60^\circ$" {font=\scriptsize}] {angle=Y--X--Z};
\end{tikzpicture}
\sgcaption{异面线所成角：平移其中一条，使两条方向线相交。}
\end{center}

\begin{examplebox}[原稿型：正方体中的异面线夹角]
  正方体 $ABCD-A_1B_1C_1D_1$ 的棱长为 $a$，求异面直线 $A_1B$ 与 $B_1C$ 所成的角。

  \textbf{最短解：}$B_1C\parallel A_1D$，故所求角为 $\angle BA_1D$。取
  $\overrightarrow{A_1B}=(a,0,-a)$、$\overrightarrow{B_1C}=(0,a,-a)$，则
  \[
    \cos\theta=\frac{a^2}{(a\sqrt2)(a\sqrt2)}=\frac12,
    \qquad \boxed{\theta=60^\circ}.
  \]
  \textbf{速查：}先平移成相交线；算角时方向向量可直接点乘。
\end{examplebox}

\section{平行关系}

\subsection{直线与平面平行}

若直线 $a$ 与平面 $\alpha$ 没有公共点，称 $a\parallel\alpha$。画图时把 $a$ 画成与平面内某条直线平行。

\begin{lawbox}[线面平行的判定与性质]
  \textbf{判定：}若 $a\not\subset\alpha$，$b\subset\alpha$ 且 $a\parallel b$，则 $a\parallel\alpha$。

  \textbf{性质：}若 $a\parallel\alpha$，过 $a$ 的平面 $\beta$ 与 $\alpha$ 交于 $b$，则 $a\parallel b$。
\end{lawbox}

判定时是在平面内“找平行线”；性质则是用过已知直线的辅助平面“截出平行线”。两者方向相反，构成最常用的线面转化：
\[
  \text{线面平行}\ \Longleftrightarrow\ \text{线线平行（配合辅助平面）}.
\]

\subsection{平面与平面平行}

若 $\alpha,\beta$ 没有公共点，称 $\alpha\parallel\beta$。

\begin{lawbox}[面面平行的判定]
  若平面 $\alpha$ 内两条相交直线 $a,b$ 都平行于平面 $\beta$，则 $\alpha\parallel\beta$。

  等价地，若 $a,b\subset\alpha$，$a',b'\subset\beta$，$a\cap b=O$，$a'\cap b'=O'$，且 $a\parallel a'$、$b\parallel b'$，则 $\alpha\parallel\beta$。
\end{lawbox}

\begin{lawbox}[面面平行的性质]
  \begin{enumerate}
    \item 若 $\alpha\parallel\beta$，则 $\alpha$ 内任一直线都平行于 $\beta$。
    \item 若 $\alpha\parallel\beta$，平面 $\gamma$ 分别与它们交于 $a,b$，则 $a\parallel b$。
    \item 夹在两个平行平面间的平行线段相等；两平行平面间的距离处处相等。
  \end{enumerate}
\end{lawbox}

\begin{center}
\begin{tikzpicture}[scale=.82]
  \path[sg plane] (0,0)--(4,0)--(4.9,1.35)--(.9,1.35)--cycle;
  \draw[very thick,BookGold] (.8,.55)--(4.05,.55) node[right,sg label] {$b$};
  \draw[very thick,BookBlue] (.55,2.1)--(3.8,2.1) node[right,sg label] {$a$};
  \draw[sg aux] (.55,2.1)--(.8,.55) (3.8,2.1)--(4.05,.55);
  \node[sg label] at (4.1,1.15) {$\alpha$};
  \node[font=\small] at (2.4,-.45) {$a\parallel b\subset\alpha\Rightarrow a\parallel\alpha$};

  \path[sg plane] (6.1,0)--(10.1,0)--(11,1.25)--(7,1.25)--cycle;
  \path[sg plane,fill=BookGold!15] (6.1,2.05)--(10.1,2.05)--(11,3.3)--(7,3.3)--cycle;
  \draw[very thick,BookBlue] (6.8,.48)--(10.25,.48);
  \draw[very thick,BookBlue] (7.35,.88)--(9.1,.15);
  \draw[very thick,BookGold] (6.8,2.53)--(10.25,2.53);
  \draw[very thick,BookGold] (7.35,2.93)--(9.1,2.2);
  \node[sg label] at (10.25,1) {$\beta$};
  \node[sg label] at (10.25,3.05) {$\alpha$};
  \node[font=\small] at (8.55,-.45) {两组相交对应线平行 $\Rightarrow\alpha\parallel\beta$};
\end{tikzpicture}
\sgcaption{判定平行的核心动作都是“在面内找线”。}
\end{center}

\section{垂直关系}

\subsection{直线与平面垂直}

若直线 $l$ 与平面 $\alpha$ 内任意直线都垂直，称 $l\perp\alpha$；交点称为垂足。

\begin{lawbox}[线面垂直的判定]
  若 $m,n\subset\alpha$，$m\cap n=O$，且 $l\perp m$、$l\perp n$，则 $l\perp\alpha$。
\end{lawbox}

“两条相交直线”不可省略：只垂直于平面内一条直线，不能推出垂直于该平面。常用推论为
\[
  l\perp\alpha,\ m\parallel l\Rightarrow m\perp\alpha;\qquad
  l\perp\alpha,\ m\perp\alpha\Rightarrow l\parallel m;\qquad
  l\perp\alpha,\ l\perp\beta\Rightarrow\alpha\parallel\beta.
\]
过空间一点有且只有一条直线垂直于已知平面；过空间一点有且只有一个平面垂直于已知直线。

\subsection{斜线、射影与三垂线定理}

设 $P\notin\alpha$，$PO\perp\alpha$。$PO$ 是垂线段，$O$ 是 $P$ 在 $\alpha$ 上的射影；对 $A\in\alpha$、$A\ne O$，$PA$ 是斜线段，$OA$ 是 $PA$ 在 $\alpha$ 内的射影。

\begin{lawbox}[三垂线定理及逆定理]
  对平面内过 $A$ 的直线 $a$：
  \[
    a\perp OA\ \Longrightarrow\ a\perp PA,
    \qquad
    a\perp PA\ \Longrightarrow\ a\perp OA.
  \]
  即：平面内一条直线垂直于斜线的射影，当且仅当它垂直于斜线本身。
\end{lawbox}

\begin{center}
\begin{tikzpicture}[scale=.9]
  \path[sg plane] (0,0)--(5.1,0)--(6.1,1.55)--(1,1.55)--cycle;
  \coordinate (O) at (2.3,.75); \coordinate (A) at (4.6,.75); \coordinate (P) at (2.3,3.35);
  \draw[sg edge] (P)--(O) (P)--(A) (O)--(A);
  \draw[very thick,BookGold] (4.6,-.1)--(4.6,1.55) node[above,sg label] {$a$};
  \draw[sg aux] (O)++(.18,0)--++(0,.18)--++(-.18,0);
  \draw[sg aux] (A)++(-.18,0)--++(0,.18)--++(.18,0);
  \node[sg label,left] at (P) {$P$}; \node[sg label,below] at (O) {$O$};
  \node[sg label,below] at (A) {$A$}; \node[sg label] at (5.3,1.25) {$\alpha$};
  \node[font=\small,align=left] at (8.4,2.2) {$PO\perp\alpha$\\$OA$ 是 $PA$ 的射影\\[3pt]
    $a\perp OA\Longleftrightarrow a\perp PA$};
\end{tikzpicture}
\sgcaption{三垂线定理把难看的空间垂直，压回平面内的射影垂直。}
\end{center}

\begin{examplebox}[原稿型：正方体的线面角]
  正方体 $ABCD-A_1B_1C_1D_1$ 棱长为 $a$，求体对角线 $AC_1$ 与底面 $ABCD$ 所成角 $\theta$。

  \textbf{最短解：}$C_1C\perp$ 底面，$AC$ 是 $AC_1$ 在底面内的射影。在直角三角形 $ACC_1$ 中，
  \[
    AC=a\sqrt2,\quad AC_1=a\sqrt3,
    \qquad \boxed{\sin\theta=\frac{CC_1}{AC_1}=\frac{\sqrt3}{3}}.
  \]
  \textbf{速查：}求线面角，先找斜线在平面内的射影，再进直角三角形。
\end{examplebox}

\subsection{二面角与面面垂直}

从一条直线出发的两个半平面组成二面角，这条直线称为棱。过棱上一点，在两个半平面内分别作垂直于棱的射线，其夹角是二面角的平面角。

若二面角的平面角为 $90^\circ$，称两平面互相垂直，记作 $\alpha\perp\beta$。

\begin{center}
\begin{tikzpicture}[scale=.9]
  \coordinate (O) at (0,0); \coordinate (L) at (5,0);
  \path[draw=BookBlue,fill=BookBlueLight] (O)--(L)--(4.1,1.9)--(.9,1.9)--cycle;
  \path[draw=BookBlue,fill=BookGold!12] (O)--(L)--(4.35,-1.6)--(.65,-1.6)--cycle;
  \coordinate (M) at (2.4,0); \coordinate (A) at (2.4,1.5); \coordinate (B) at (2.4,-1.25);
  \draw[very thick,BookGold] (A)--(M)--(B);
  \pic[draw=BookGold,angle radius=7mm,"$\theta$" {font=\small}] {angle=A--M--B};
  \node[sg label,above left] at (L) {$l$};
  \node[sg label] at (4.05,1.35) {$\alpha$}; \node[sg label] at (4,-1.2) {$\beta$};
  \node[font=\small,align=left] at (8.05,.25) {$AM\perp l,\ BM\perp l$\\$\angle AMB$ 是平面角};
\end{tikzpicture}
\sgcaption{求二面角：在棱上选一点，分别在两面内作棱的垂线。}
\end{center}

\begin{lawbox}[面面垂直的判定与性质]
  \textbf{判定：}若平面 $\beta$ 经过平面 $\alpha$ 的一条垂线，则 $\alpha\perp\beta$。

  \textbf{性质：}若 $\alpha\perp\beta$，$\alpha\cap\beta=l$，直线 $a\subset\alpha$ 且 $a\perp l$，则 $a\perp\beta$。
\end{lawbox}

\subsection{空间角的两个高密度公式}

若直线与三条两两垂直的坐标轴所成角分别为 $\alpha,\beta,\gamma$，则方向余弦满足
\[
  \cos^2\alpha+\cos^2\beta+\cos^2\gamma=1,
  \qquad
  \sin^2\alpha+\sin^2\beta+\sin^2\gamma=2.
\]

\section{多面体}

由若干平面多边形围成的几何体称为多面体；围成它的多边形是面，相邻面的公共边是棱，棱的公共端点是顶点。若多面体位于每个面的同一侧，称为凸多面体，否则为凹多面体。

\subsection{棱柱、棱锥与棱台}

\begin{description}
  \item[棱柱] 两个面互相平行，其余各面为平行四边形，且相邻侧面的公共边互相平行。平行的两个面是底面，其余为侧面。侧棱垂直底面的棱柱是直棱柱；底面为正多边形的直棱柱是正棱柱。
  \item[平行六面体] 底面为平行四边形的四棱柱。六个面都是矩形时为长方体；六个面都是正方形时为正方体。
  \item[棱锥] 一个面是多边形，其余各面是有公共顶点的三角形。若底面为正多边形且顶点在底面的射影为底面中心，称为正棱锥。
  \item[棱台] 用平行于棱锥底面的平面截棱锥，截面与原底面之间的部分称为棱台；由正棱锥截得的是正棱台。
\end{description}

\begin{center}
\begin{tikzpicture}[scale=.72]
  % prism
  \draw[sg edge] (0,0) rectangle (2.1,2.5);
  \draw[sg edge] (.75,.65) rectangle (2.85,3.15);
  \draw[sg edge] (0,0)--(.75,.65) (2.1,0)--(2.85,.65) (0,2.5)--(.75,3.15) (2.1,2.5)--(2.85,3.15);
  \draw[sg hidden] (.75,.65)--(2.85,.65);
  \node[font=\small] at (1.4,-.55) {棱柱};
  % pyramid
  \coordinate (S) at (5.5,3.35); \coordinate (A) at (3.9,0); \coordinate (B) at (6.3,0); \coordinate (C) at (7.15,.8); \coordinate (D) at (4.75,.8);
  \draw[sg edge] (A)--(B)--(C) (A)--(S)--(B) (S)--(C);
  \draw[sg hidden] (C)--(D)--(A) (D)--(S);
  \draw[sg aux] (S)--(5.5,.42);
  \node[font=\small] at (5.45,-.55) {棱锥};
  % frustum
  \draw[sg edge] (8.5,0)--(11,0)--(11.7,.65) (8.5,0)--(9.2,.65);
  \draw[sg hidden] (9.2,.65)--(11.7,.65);
  \draw[sg edge] (9.25,2.45)--(10.65,2.45)--(11.05,2.8)--(9.65,2.8)--cycle;
  \draw[sg edge] (8.5,0)--(9.25,2.45) (11,0)--(10.65,2.45) (11.7,.65)--(11.05,2.8);
  \draw[sg hidden] (9.2,.65)--(9.65,2.8);
  \node[font=\small] at (10.05,-.55) {棱台};
\end{tikzpicture}
\sgcaption{先识别“两个平行底面 / 一个公共顶点 / 平行截去尖端”，再选公式。}
\end{center}

正棱锥的侧棱相等，各侧面是全等的等腰三角形；侧面等腰三角形底边上的高称为斜高。设底面周长为 $p$、斜高为 $l$，则
\[
  S_{\text{正棱锥侧}}=\frac12pl.
\]
正棱台上下底周长为 $p,p'$、斜高为 $l$，则
\[
  S_{\text{正棱台侧}}=\frac12(p+p')l.
\]

\begin{examplebox}[原稿型：正棱锥侧面积]
  正四棱锥底面边长为 $a$，斜高为 $l$，求侧面积；若高为 $h$，求体积。

  \textbf{最短解：}四个侧面是全等三角形，底面周长 $p=4a$，故
  \[
    \boxed{S_{\text{侧}}=\frac12pl=2al},
    \qquad
    \boxed{V=\frac13a^2h}.
  \]
  \textbf{速查：}表面积沿面展开；体积只认“底面积 $\times$ 高”。斜高 $l$ 不是体高 $h$。
\end{examplebox}

\section{旋转体、表面积与体积}

矩形、直角三角形、直角梯形分别绕其一条直角边旋转，得到圆柱、圆锥、圆台；旋转轴所在的截面称为轴截面。

\begin{center}
\begin{tikzpicture}[scale=.72]
  \draw[sg edge] (0,0) ellipse (1.2 and .35) (0,2.7) ellipse (1.2 and .35);
  \draw[sg edge] (-1.2,0)--(-1.2,2.7) (1.2,0)--(1.2,2.7);
  \draw[sg hidden] (-1.2,0) arc (180:360:1.2 and .35) (-1.2,2.7) arc (180:360:1.2 and .35);
  \draw[sg aux] (0,-.35)--(0,3.05); \node[font=\small] at (0,-.8) {圆柱};
  \draw[sg edge] (4.2,0) ellipse (1.25 and .35) (2.95,0)--(4.2,3.05)--(5.45,0);
  \draw[sg hidden] (2.95,0) arc (180:360:1.25 and .35);
  \draw[sg aux] (4.2,-.35)--(4.2,3.05); \node[font=\small] at (4.2,-.8) {圆锥};
  \draw[sg edge] (8.5,0) ellipse (1.35 and .38) (8.5,2.65) ellipse (.72 and .23);
  \draw[sg edge] (7.15,0)--(7.78,2.65) (9.85,0)--(9.22,2.65);
  \draw[sg hidden] (7.15,0) arc (180:360:1.35 and .38) (7.78,2.65) arc (180:360:.72 and .23);
  \draw[sg aux] (8.5,-.38)--(8.5,2.9); \node[font=\small] at (8.5,-.8) {圆台};
\end{tikzpicture}
\sgcaption{轴截面分别是矩形、等腰三角形、等腰梯形；母线 $l$ 在轴截面里。}
\end{center}

\begin{center}
  \begin{tabular}{llll}
    \toprule
    几何体 & 侧面积 & 全面积 & 体积 \\
    \midrule
    直棱柱 & $ph$ & $ph+2S$ & $Sh$ \\
    圆柱 & $2\pi rl$ & $2\pi r(r+l)$ & $\pi r^2h$ \\
    圆锥 & $\pi rl$ & $\pi r(r+l)$ & $\frac13\pi r^2h$ \\
    圆台 & $\pi(r_1+r_2)l$ & $\pi(r_1+r_2)l+\pi(r_1^2+r_2^2)$ & $\frac13\pi h(r_1^2+r_1r_2+r_2^2)$ \\
    \bottomrule
  \end{tabular}
\end{center}

其中 $S$ 为底面积，$p$ 为底面周长，$h$ 为高，$l$ 为母线或斜高。

\subsection{祖氏原理与统一体积公式}

\begin{lawbox}[祖氏原理（等幂等积）]
  夹在两个平行平面之间的两个几何体，若在任意等高处的截面面积都相等，则两几何体体积相等。
\end{lawbox}

它把柱体与柱体、锥体与锥体的体积统一起来：
\[
  V_{\text{柱体}}=Sh,
  \qquad
  V_{\text{锥体}}=\frac13Sh,
  \qquad
  V_{\text{台体}}=\frac h3\left(S+\sqrt{SS'}+S'\right).
\]

台体公式来自补成大锥体。若小锥与大锥相似，底面积分别为 $S',S$，对应高为 $x,x+h$，则
\[
  \frac{x}{x+h}=\sqrt{\frac{S'}S};
\]
用“大锥体积减小锥体积”消去 $x$ 即得台体公式。体积题的基本路线只有三类：直接套公式、等体积转化、分割与补形。

\begin{examplebox}[原稿方法：四面体换底]
  四面体体积既可把 $\triangle ABC$ 作为底，也可把 $\triangle PAB$ 作为底：
  \[
    V_{P-ABC}=\frac13S_{ABC}\,d(P,ABC)
    =\frac13S_{PAB}\,d(C,PAB).
  \]
  \textbf{速查：}题目给哪一个面面积、哪一条垂距，就把那个面当底；“换底”等体积常比硬求空间高更快。
\end{examplebox}

\section{球}

半圆绕直径旋转所得的几何体是球；半圆弧旋转所得曲面是球面。球心记为 $O$，半径为 $R$，直径为 $2R$。

球面被任一平面截得的截线是圆。若球心到截面的距离为 $d$、截面圆半径为 $r$，则
\[
  r^2+d^2=R^2.
\]

\begin{center}
\begin{tikzpicture}[scale=.9]
  \draw[sg edge] (0,0) circle (2.05);
  \draw[sg edge] (0,0) ellipse (2.05 and .62);
  \draw[sg hidden] (-2.05,0) arc (180:360:2.05 and .62);
  \coordinate (O) at (0,0); \coordinate (H) at (0,.62); \coordinate (A) at (1.72,.33);
  \draw[very thick,BookGold] (O)--(H)--(A) (O)--(A);
  \draw[sg aux] (H)++(.15,-.03)--++(-.03,-.15)--++(-.15,.03);
  \node[sg label,left] at (O) {$O$}; \node[sg label,above] at (H) {$H$}; \node[sg label,right] at (A) {$A$};
  \node[sg label,left] at ($(O)!.5!(H)$) {$d$};
  \node[sg label,above] at ($(H)!.5!(A)$) {$r$};
  \node[sg label,below right] at ($(O)!.5!(A)$) {$R$};
  \node[font=\small,align=left] at (5.3,.5) {球题的核心截面：\\[3pt]$OH\perp$ 截面，$HA=r$，$OA=R$\\[3pt]
    $\boxed{R^2=d^2+r^2}$};
\end{tikzpicture}
\sgcaption{先把球切成一个过球心的平面图，再用勾股定理。}
\end{center}
过球心的截面圆称为大圆；球面上两点间的最短球面距离沿过两点的大圆劣弧量取。

\[
  S_{\text{球}}=4\pi R^2,
  \qquad
  V_{\text{球}}=\frac43\pi R^3.
\]

\subsection{内切球与外接球}

球面与多面体各面相切时，球为多面体的内切球。设内切球半径为 $r$、多面体表面积为 $S_{\text{表}}$，把球心与各面连接并分割成若干棱锥，有
\[
  V=\frac13rS_{\text{表}}.
\]

球面经过多面体全部顶点时，球为多面体的外接球。正四面体棱长为 $a$ 时：
\[
  h=\frac{\sqrt6}{3}a,
  \qquad
  r=\frac{\sqrt6}{12}a,
  \qquad
  R=\frac{\sqrt6}{4}a,
  \qquad
  R=3r.
\]

处理内外接球题，先找对称轴与截面：把空间问题压缩为包含球心、切点或顶点的平面图形，再用勾股定理。

\begin{examplebox}[原稿例：截面圆反求球半径]
  球面上一截面内有 $\triangle ABC$，$AB=BC=2$、$AC=2\sqrt2$；球心到截面距离为 $d=R/2$。求球的表面积与体积。

  \textbf{最短解：}$AB^2+BC^2=AC^2$，故 $\triangle ABC$ 为直角三角形，截面圆半径等于斜边一半：$r=AC/2=\sqrt2$。于是
  \[
    R^2=r^2+d^2=2+\frac{R^2}{4}
    \Rightarrow R^2=\frac83,
  \]
  \[
    \boxed{S=\frac{32\pi}{3}},
    \qquad
    \boxed{V=\frac{64\sqrt6\pi}{27}}.
  \]
  \textbf{速查：}截面内先求外接圆半径 $r$，再用 $R^2=r^2+d^2$ 升回空间。
\end{examplebox}

\begin{examplebox}[原稿例：正四面体的内、外接球]
  正四面体棱长为 $a$。底面中心记为 $O$，顶点为 $P$。底面高为 $\sqrt3a/2$，故顶点到底面中心的水平距离为 $a/\sqrt3$；由勾股定理
  \[
    PO=h=\sqrt{a^2-\frac{a^2}{3}}=\frac{\sqrt6}{3}a.
  \]
  内、外球心都在 $PO$ 上；利用体积分割或相似三角形可得
  \[
    \boxed{r=\frac h4=\frac{\sqrt6}{12}a},
    \qquad
    \boxed{R=\frac{3h}{4}=\frac{\sqrt6}{4}a},
    \qquad R=3r.
  \]
\end{examplebox}

\section{投影、三视图与直观图}

光线照射物体，在平面上形成的影子是投影。投影线交于一点的是中心投影；投影线互相平行的是平行投影。平行投影又分为斜投影和正投影（投影线垂直投影面）。

平行投影保持共线性和平行性；与投影面平行的线段、平面图形保持真实长度与形状。它一般不保持角度，也不保持不与投影面平行的长度。

\subsection{三视图}

正视图、侧视图、俯视图分别反映物体的
\[
  \text{长与高},\qquad\text{宽与高},\qquad\text{长与宽}.
\]
排图遵循“长对正、高平齐、宽相等”。可见轮廓画实线，不可见棱画虚线。由三视图还原几何体时，先由俯视图定底面位置，再由正、侧视图补高度。

\begin{center}
\begin{tikzpicture}[scale=.78]
  % object
  \draw[sg edge] (0,0) rectangle (2.3,1.7);
  \draw[sg edge] (.65,.55) rectangle (2.95,2.25);
  \draw[sg edge] (0,0)--(.65,.55) (2.3,0)--(2.95,.55) (0,1.7)--(.65,2.25) (2.3,1.7)--(2.95,2.25);
  \draw[sg hidden] (.65,.55)--(2.95,.55);
  \node[font=\small] at (1.45,-.5) {物体};
  \draw[-{Latex[length=2mm]},BookGold,thick] (3.35,1.1)--(4.25,1.1);
  % front
  \draw[sg edge] (4.7,.25) rectangle (6.65,1.75);
  \node[font=\small] at (5.68,-.25) {正视：长 $\times$ 高};
  % side
  \draw[sg edge] (8.1,.25) rectangle (9.45,1.75);
  \node[font=\small] at (8.78,-.25) {侧视：宽 $\times$ 高};
  % top
  \draw[sg edge] (4.7,2.75) rectangle (6.65,3.85);
  \node[font=\small] at (5.68,4.2) {俯视：长 $\times$ 宽};
  \draw[sg aux,<->] (4.7,2.42)--(6.65,2.42) node[midway,below,sg label] {长对正};
  \draw[sg aux,<->] (7.1,.25)--(7.1,1.75) node[midway,rotate=90,above,sg label] {高平齐};
\end{tikzpicture}
\sgcaption{三视图不是三张独立的图：同一尺寸必须在相邻视图中对齐。}
\end{center}

\begin{examplebox}[原稿型：由三视图还原]
  若俯视图是 $2\times2$ 方格，四格均有物体；正视图两列最高层数为 $2,1$，侧视图两列最高层数也为 $2,1$，则先在俯视图四格放一层，再把唯一的第二层放在同时满足两张“高度为 2”投影的交叉格。

  \textbf{速查：}俯视图管“哪里有”；正、侧视图管“每一列最高多高”。还原题本质是两组最大值约束的交叉定位。
\end{examplebox}

\subsection{斜二测直观图}

通常取 $x$ 轴水平、$z$ 轴竖直、$y$ 轴与 $x$ 轴成 $45^\circ$ 或 $135^\circ$。与 $x,z$ 轴平行的长度不变，与 $y$ 轴平行的长度取原来的一半；平行关系保持不变。先画底面，再沿 $z$ 轴补高度，最后连接并区分实、虚线。

\section{三余弦定理}

设 $P\notin\alpha$，$PO\perp\alpha$，$A,B\in\alpha$。$OA$ 是斜线 $PA$ 在平面 $\alpha$ 内的射影。令
\[
  \angle PAO=\theta,
  \qquad
  \angle OAB=\varphi,
  \qquad
  \angle PAB=\psi.
\]

\begin{lawbox}[三余弦定理]
  \[
    \boxed{\cos\psi=\cos\theta\cos\varphi}.
  \]
  即：斜线与平面内直线夹角的余弦，等于斜线与其射影夹角的余弦乘以射影与该直线夹角的余弦。
\end{lawbox}

\begin{center}
\begin{tikzpicture}[scale=.92]
  \path[sg plane] (0,0)--(6.1,0)--(7.1,1.75)--(1,1.75)--cycle;
  \coordinate (O) at (2.35,.8); \coordinate (A) at (4.6,.8); \coordinate (B) at (5.75,1.45); \coordinate (P) at (2.35,3.6);
  \draw[sg edge] (P)--(O)--(A)--(P) (A)--(B);
  \draw[sg hidden] (O)--(B);
  \draw[sg aux] (O)++(.17,0)--++(0,.17)--++(-.17,0);
  \pic[draw=BookGold,angle radius=6mm,"$\theta$" {font=\scriptsize}] {angle=P--A--O};
  \pic[draw=BookBlue,angle radius=8mm,"$\varphi$" {font=\scriptsize}] {angle=O--A--B};
  \pic[draw=BookGold,angle radius=11mm,"$\psi$" {font=\scriptsize}] {angle=P--A--B};
  \foreach \p/\pos in {P/above,O/below,A/below,B/right}\node[sg label,\pos] at (\p) {$\p$};
  \node[font=\small,align=left] at (9.2,1.75) {$PA$ 的射影是 $OA$\\[3pt]
    $\theta=\angle PAO$\\$\varphi=\angle OAB$\\$\psi=\angle PAB$\\[4pt]
    $\boxed{\cos\psi=\cos\theta\cos\varphi}$};
\end{tikzpicture}
\sgcaption{三余弦定理同时追踪“斜线—射影—面内目标线”三个角。}
\end{center}

证明只用正射影。因为 $PO\perp AB$，
\[
  \overrightarrow{AP}\cdot\overrightarrow{AB}
  =\left(\overrightarrow{AO}+\overrightarrow{OP}\right)\cdot\overrightarrow{AB}
  =\overrightarrow{AO}\cdot\overrightarrow{AB}.
\]
两边除以 $|AP|\,|AB|$，再插入 $|AO|/|AP|=\cos\theta$，即得结论。令 $\psi=90^\circ$，便得到三垂线定理及其逆定理。

\section{判定链速查}

\begin{center}
  \begin{tabularx}{\textwidth}{@{}lX@{}}
    \toprule
    目标 & 最短判定链 \\
    \midrule
    线在面内 & 线上两点在面内 \\
    点共线 & 同时属于两个相交平面 \\
    线面平行 & 面外直线 $\parallel$ 面内直线 \\
    面面平行 & 一面内两相交直线分别 $\parallel$ 另一面 \\
    线面垂直 & 直线垂直于面内两相交直线 \\
    面面垂直 & 一平面经过另一平面的垂线 \\
    异面线垂直 & 平移成相交直线后夹角为 $90^\circ$ \\
    斜线垂直面内线 & 转化为其射影垂直该线 \\
    \bottomrule
  \end{tabularx}
\end{center}
